% =============================================================================
\lecture[number=12, short=答题方法与答案]{答题方法与练一练答案}
% =============================================================================

\begin{keypoints}
\begin{itemize}
  \item 计算题：先画图，再列式；写清设定和单位；用“楔子 $x+y=t$”“剩余加起来不变”等关系检查结果。
  \item 分析题四步：\textbf{结论 $\to$ 机制 $\to$ 条件 $\to$ 另一面}。用模型说明“为什么”，再说明模型在什么条件下成立、现实中还要考虑什么。
  \item 社科题先找“关键词”，想到对应的工具：比率想到分母和权重，调查想到偏差，“有关”想到混杂和随机分组。
\end{itemize}
\end{keypoints}

\section{计算题}

\begin{enumerate}
  \item \textbf{先画图}。哪怕是草图，也能帮你看清哪条线动、面积是哪一块。经济题里很多“面积”就是三角形和矩形。
  \item \textbf{写清设定}。例如“设买方价格为 $P_{\text{买}}$，卖方到手价格为 $P_{\text{买}}-t$”；单位（元/件、月利率还是年利率）写出来。
  \item \textbf{检查}。几个随手可用的检查：买方多付的加卖方少得的等于税额；各块面积加起来等于原来的总剩余；贴现后的还款总现值等于借款额；
        两两表决的票数加起来等于总人数。
\end{enumerate}

\section{分析题怎么写}

分析题、论述题最怕“只有结论没有道理”或“只有一堆名词”。一个稳妥的结构是四步（图~\ref{fig:es-steps}）。

\begin{figure}[!htbp]
\centering
\begin{tikzpicture}[st/.style={node box, minimum height=1.35cm, minimum width=2.9cm},
    fl/.style={-{Stealth[length=2.4mm,width=1.8mm]}, draw=mjInk!70, line width=.8pt}]
  \node[st, draw=mjAccent, fill=mjAccent!8] (a) at (0,0) {\textbf{1. 结论}\\ 一句话回答问题};
  \node[st, draw=mjPine, fill=mjPine!8] (b) at (3.55,0) {\textbf{2. 机制}\\ 用哪个模型、\\哪条线怎样动};
  \node[st, draw=mjOchre, fill=mjOchre!10] (c) at (7.1,0) {\textbf{3. 条件}\\ 模型的假设、\\成立的前提};
  \node[st, draw=mjSecond, fill=mjSecond!8] (d) at (10.65,0) {\textbf{4. 另一面}\\ 现实的复杂性、\\局限与建议};
  \draw[fl] (a) -- (b); \draw[fl] (b) -- (c); \draw[fl] (c) -- (d);
\end{tikzpicture}
\caption{分析题的四步结构}
\label{fig:es-steps}
\end{figure}

\begin{example}[一个示范]
\textbf{问}：为什么节假日的热门车票总是“一票难求”？

\textbf{答}：（结论）因为票价被固定在低于高峰期均衡价格的水平，形成了短缺。
（机制）节假日出行需求大幅增加，需求曲线右移，而运力短期内难以增加，供给曲线接近竖直；票价不随之上涨，想买的人数远多于座位数（第 4 讲最高限价）。
价格不再承担分配的功能，只好靠排队、抢票来分配，抢票的时间和精力就是隐性的代价。
（条件）这一分析假设票价不能灵活调整、运力短期固定。
（另一面）完全按供求涨价能减少短缺，但会让收入较低、又必须回家的人承受很高的票价，所以现实中常保持票价稳定，同时增开临时列车、实行实名制和分时段放票，在效率和公平之间做权衡。
\end{example}

\section{看到什么，想到什么}

\begin{center}
\small
\begin{tabular}{p{4.2cm}p{8.8cm}}
  \toprule
  题目里出现 & 想到 \\
  \midrule
  价格、征税、补贴、限价 & 供求图；哪条线动；弹性；剩余与无谓损失（第 1--4 讲） \\
  污染、拥堵、疫苗 & 外部性；社会成本；庇古税、排污交易（第 4 讲） \\
  两家企业、合作与背叛 & 收益表、纳什均衡、囚徒困境；重复博弈与贴现因子（第 5 讲） \\
  利息、分期、贷款、“哪个划算” & 现值；年金公式；同一时刻比较（第 6 讲） \\
  分工、贸易、“谁更强” & 机会成本、比较优势；交换比例范围；分配效应（第 7 讲） \\
  投票、表决、规则 & 孔多塞循环、议程、不同规则不同赢家、策略性投票、阿罗定理、中位选民（第 8 讲） \\
  比率、百分比、“合计” & 分母；加权平均；辛普森悖论；基础比率（第 9 讲） \\
  调查、民调、问卷 & 抽样框、无回答、措辞；随机误差与偏差；概率抽样（第 10 讲） \\
  “A 和 B 有关”“因此应该……” & 混杂、反向因果、自我选择；随机对照实验；双重差分（第 11 讲） \\
  \bottomrule
\end{tabular}
\end{center}

\begin{pitfall}
\begin{itemize}
  \item 名词堆砌不得分：写了“阿罗不可能定理”，却没有说明它和本题的关系，等于没写。
  \item 开放题要有\textbf{具体}的做法：“扩大样本”“加强监管”这类笼统的话，换成“按年级分层、随机抽取班级、匿名填写一周的睡眠日志”这样具体的步骤。
  \item 题目给了定义（如“纳什均衡指……”），就严格按题目的定义来做，不要套用自己记得的其他说法。
\end{itemize}
\end{pitfall}

\section{练一练答案}

\paragraph{第 1 讲}
\begin{enumerate}[label=1.\arabic*, leftmargin=2.6em]
  \item $120-4P=2P$，均衡 $P=20$，$Q=40$。价格定为 15 元时，需求量 60，供给量 30，\textbf{短缺} 30。
  \item (a) 电池是电动车的主要成本，成本下降使\textbf{供给}右移：价格下降，数量增加。
        (b) 汽油涨价使燃油车用起来更贵，电动车作为替代品的\textbf{需求}右移：价格上升，数量增加。
        (c) 两者都使数量增加，所以数量一定增加；一个使价格下降、一个使价格上升，价格变化不确定。
  \item 观察到的“价格上涨、成交量增加”，是不同时期的供求曲线的不同交点，不是同一条需求曲线上的点。很可能是收入增长、人口流入、对房价继续上涨的预期等使需求曲线右移了，
        新交点的价格和数量都更高。需求曲线向下倾斜说的是“其他条件不变时”价格与需求量的关系，这个现象并不与它矛盾。
\end{enumerate}

\paragraph{第 2 讲}
\begin{enumerate}[label=2.\arabic*, leftmargin=2.6em]
  \item 需求曲线与纵轴交于 $P=30$。均衡 $P=20$，$Q=40$。消费者剩余 $\frac12\times40\times(30-20)=200$，生产者剩余 $\frac12\times40\times20=400$。
  \item 价格上涨 $20\%$，游客减少 $10\%$，$|E|\approx\frac{10\%}{20\%}=0.5$，缺乏弹性。门票收入从 $200000$ 元增加到 $216000$ 元。
  \item 临时订票的多是商务出行或有急事的人，没有别的选择，需求缺乏弹性，航空公司定高价也不会少卖多少；提前订票的多是可去可不去的旅游者，需求富有弹性，降价能吸引更多人。
        对弹性不同的人收不同的价格（\textbf{价格歧视}），能增加航空公司的收入。
\end{enumerate}

\paragraph{第 3 讲}
\begin{enumerate}[label=3.\arabic*, leftmargin=2.6em]
  \item $b=4$，$d=2$，$t=3$。买方价格上升 $x=\frac{2}{6}\times3=1$，卖方价格下降 $y=\frac46\times3=2$。原均衡 $P=20$，所以买方共付 21 元（其中标价 18 元，另交税 3 元），卖方到手 18 元，
        成交量 $2\times18=36$（检验：$120-4\times21=36$）。税收 $3\times36=108$ 元。买方每件负担 1 元，卖方负担 2 元。无谓损失 $\frac12\times3\times(40-36)=6$。
  \item 记 $k=\frac{bd}{b+d}$，则无谓损失 $=\frac12kt^2$，$t$ 加倍后变为 $\frac12k(2t)^2$，是原来的 4 倍。税收 $R(t)=t(Q^*-kt)$，
        $R(2t)=2t(Q^*-2kt)=2R(t)-2kt^2<2R(t)$。
  \item 富人购买豪华跑车，有很多替代选择（别的奢侈品、去国外买、推迟购买），需求的弹性可能很大，那么税负会有相当一部分落在生产商、经销商和相关工人身上，未必主要由富人承担。
        要判断，还需要知道需求和供给两边的弹性：例如短期内工厂产能固定、供给缺乏弹性，卖方负担会更多。
\end{enumerate}

\paragraph{第 4 讲}
\begin{enumerate}[label=4.\arabic*, leftmargin=2.6em]
  \item (1) 价格 25 元时，需求量 20，供给量 50，过剩 30。(2) $30\times25=750$ 元。
  \item (1) 工厂只考虑自己的成本，不考虑渔民的损失，每吨的社会成本比工厂的成本高 50 元，所以生产到了对社会来说不值得的数量。
        (2) 每吨征税 50 元后，工厂自己的成本等于社会成本，它会自己把产量减到对社会最好的水平。
        (3) 工厂每生产 1 吨要赔 50 元，效果和征税 50 元一样，产量同样减到有效率的水平；区别是这 50 元给了受损的渔民而不是政府。反过来，即使法律允许工厂排污，只要协商成本低，渔民也会出钱请工厂减产（每少产 1 吨，渔民最多愿出 50 元），产量同样会达到有效率的水平；
        产权给谁只影响钱归谁。这就是科斯的思想：产权明确、协商成本低时，私人之间就能解决外部性。
  \item 偏少。个人决定是否接种时只考虑自己得到的保护，不考虑对周围人的保护（正外部性），所以接种的人数会少于对社会最好的人数。这是政府补贴或免费提供疫苗的理由。
\end{enumerate}

\paragraph{第 5 讲}
\begin{enumerate}[label=5.\arabic*, leftmargin=2.6em]
  \item 收益表中，（猎鹿，猎鹿）和（抓兔，抓兔）都是纳什均衡；另两格中，猎鹿的一方得 0，会改去抓兔。（猎鹿，猎鹿）对双方都更好；
        但抓兔更“保险”：不论对方怎么做，抓兔都能得 3，猎鹿则可能因为对方不合作而一无所获。所以对对方没有把握时，人们可能停在较差的均衡上。
  \item (1) $\delta\ge\frac{6-4}{6-2}=\frac12$。(2) 继续的概率 0.8 起到贴现因子的作用，$\delta=0.8\ge\frac12$，合作可以维持。
  \item 在位者打价格战得 6，容忍得 5，会打价格战；进入者预见到这一点，进入得 $-2$，不进入得 0，选择不进入。结果是（不进入）。
        扩建产能改变了在位者事后的利益，使“打价格战”的威胁变得可信，从而把潜在的竞争者挡在门外。
\end{enumerate}

\paragraph{第 6 讲}
\begin{enumerate}[label=6.\arabic*, leftmargin=2.6em]
  \item 4 年后的 10000 元，现值为 $\frac{10000}{1.04^4}\approx8548$ 元，比现在领的 8500 元略多，4 年后领更值钱（前提是利率确实是 4\%，且奖学金一定能兑现）。
  \item 记 $v=\frac1{1+i}$，现值 $=Av+Av^2+\cdots+Av^n=Av\cdot\frac{1-v^n}{1-v}$。因为 $\frac{v}{1-v}=\frac{1/(1+i)}{i/(1+i)}=\frac1i$，所以现值 $=A\cdot\frac{1-(1+i)^{-n}}{i}$。
        $n\to\infty$ 时 $(1+i)^{-n}\to0$，现值趋于 $\frac Ai$。
  \item 每月还本金 1000 元，第 $k$ 个月利息为 $(12000-1000(k-1))\times1\%$，总利息 $=120\times\frac{12+1}{2}=780$ 元。第 1 个月还款 $1000+120=1120$ 元。
  \item $72\div3=24$ 年。精确值 $\frac{\ln2}{\ln1.03}\approx23.4$ 年。
\end{enumerate}

\paragraph{第 7 讲}
\begin{enumerate}[label=7.\arabic*, leftmargin=2.6em]
  \item 小王写 1 页作文要放弃 2 道数学题，小李只要放弃 1 道。小李写作文的机会成本低，小王做数学题的机会成本低（$\frac12$ 页作文，小李是 1 页）。
        所以小王做数学题、小李写作文，虽然小王两样都比小李快。
  \item 不会。1 单位粮换 2.5 单位布，比甲国自己种粮的机会成本（2 单位布）还高，甲国宁可自己种粮，不会买粮；只有乙国愿意卖。交换比例必须在 1 和 2 之间。
  \item 工资低反映的是劳动生产率低，它不可能在所有产品上都有比较优势：两国的机会成本不同，就存在双方都获益的分工。它在劳动密集型产品上有比较优势，
        我们就在其他产品上有比较优势，通过贸易双方都能消费得更多。但国内与进口竞争的行业会受损，需要转岗培训、社会保障等政策来帮助受损者。
\end{enumerate}

\paragraph{第 8 讲}
\begin{enumerate}[label=8.\arabic*, leftmargin=2.6em]
  \item 相对多数：A 4 票，B 3 票，C 2 票，\textbf{A 胜}。两轮决选：A 与 B 进入第二轮，C 的 2 个支持者更喜欢 B，B 以 5 比 4 胜，\textbf{B 胜}。
        波达计数：A $4\times2=8$ 分，B $4+3\times2+2=12$ 分，C $3+2\times2=7$ 分，\textbf{B 胜}。两两表决：B 对 A 为 5 比 4，B 对 C 为 7 比 2，\textbf{B 是孔多塞赢家}。
  \item 例：3 人排序 A，B，C；2 人排序 B，C，A。两两表决 A 对 B、A 对 C 都是 3 比 2，A 是孔多塞赢家；但波达计数 A 得 6 分，B 得 $3+4=7$ 分，B 胜。
  \item 中位数是 5 公里，选 5 公里。加入两位后，理想位置为 2、3、5、9、10、10、10，中位数变为 9 公里，选 9 公里。
\end{enumerate}

\paragraph{第 9 讲}
\begin{enumerate}[label=9.\arabic*, leftmargin=2.6em]
  \item 小结石：A $\frac{81}{87}\approx93\%$，B $\frac{234}{270}\approx87\%$；大结石：A $\frac{192}{263}\approx73\%$，B $\frac{55}{80}\approx69\%$；
        全部：A $\frac{273}{350}=78\%$，B $\frac{289}{350}\approx83\%$。A 在每组都更好，合计却更差：A 法多用于较难治的大结石（263 人），B 法多用于小结石（270 人）。
        结石大小同时影响“用哪种方法”和“治愈率”，应当分组比较。
  \item (1) 设 100000 人：100 人有病，99 人阳性；99900 人没病，999 人阳性。阳性者有病的概率 $\frac{99}{1098}\approx9.0\%$。
        (2) 第二次仍阳性：有病的 99 人中约 98 人，没病的 999 人中约 10 人，概率 $\approx\frac{98.01}{98.01+9.99}\approx90.7\%$。
  \item 可以追问：中位数是多少？包括管理层吗？是税前还是税后，含不含奖金？我应聘的这个岗位的薪酬范围是多少？统计的是哪一年、多少人？
\end{enumerate}

\paragraph{第 10 讲}
\begin{enumerate}[label=10.\arabic*, leftmargin=2.6em]
  \item 只有打开 App 的人才会看到弹窗，不满意而不再使用的人已经离开了（覆盖偏差／幸存者偏差）；很多人直接关闭弹窗，回答的是自愿回应的一部分人（无回答偏差）；
        在 App 里被 App 自己询问，人们可能顺手点“满意”，选项的设计也可能有诱导。
  \item 误差与 $\frac1{\sqrt n}$ 成正比，误差减半，样本量要变为 4 倍。
  \item $P(\text{答“是”})=\frac46p+\frac16=0.4$，$p=(0.4-\frac16)\times\frac32=0.35$。
  \item 问题：只调查了住校生；晚自习后的这个时间、这个地点出现的人不能代表全体；“随机拦人”实际上由访问者挑选，不是概率抽样；一次口头回忆误差大。
        改进：用全校名册按年级分层，随机抽取班级（或学生）；匿名填写，连续记录一周的睡眠日志；追访没有交回的人；按各年级人数加权估计全校平均。
\end{enumerate}

\paragraph{第 11 讲}
\begin{enumerate}[label=11.\arabic*, leftmargin=2.6em]
  \item 可能的问题：家庭条件好、作息规律、自律的学生更可能每天吃早餐，成绩也更好（共同原因、自我选择）；即使吃早餐有用，强制执行的效果也可能不同。
        方案：在学校随机选择一部分班级提供免费早餐并督促食用，其他班级作对照，比较一段时间前后两组成绩的变化。
  \item 效果 $=(32\%-30\%)-(33\%-28\%)=-3$ 个百分点，即这项措施使近视率的上升少了约 3 个百分点。假设：如果不推行这项措施，两类学校近视率的变化会相同（平行趋势），
        并且同一时期没有其他只影响试点学校的变化。
  \item 样本是随机抽取的，所以“参加者的成绩提高更多”这一现象可以推广到全国高中生；但是否参加是学生自己选择的，参加者可能本来就更有动力，
        所以不能说明成绩提高是训练营造成的。要判断效果，应当在报名者中随机决定谁参加。
\end{enumerate}
